\documentclass[12pt]{article}
\usepackage[utf8]{inputenc}
\usepackage[T1]{fontenc}
\usepackage{hyperref}
\title{Sikaritai Language Resources for K-12 Curriculum Development}
\author{Compiled by Claude Sonnet 5 (claude-sonnet-5) for Andr\'as Kornai}
\date{}
\begin{document}
\maketitle

\section{Introduction}
Sikaritai (ISO 639-3 tty; Glottocode sika1263) is a Lakes Plain language of Papua, Indonesia. It is named after Sikari village in Rufaer District, Mamberamo Raya Regency, and is spoken near where the Idenburg and Rouffaer Rivers meet the Mamberamo River. Wikipedia gives about 2,000 speakers (2015) and Joshua Project gives 1,200 (citing Ethnologue 2016). Sikaritai, Obokuitai and Eritai form a dialect cluster. Little material was found.

\section{Primary Educational Resources}
None were found.

\section{Linguistic and Reference Materials}
\begin{itemize}
\item Sikaritai language (Wikipedia): \url{https://en.wikipedia.org/wiki/Sikaritai_language}
\item Clouse (1997), a source on Lakes Plain classification and lexicon, with page images on the Trans New Guinea database: \url{https://transnewguinea.org/source/clouse-1997?page=109}
\item Lakes Plain languages (Wikipedia mirror): \url{https://static.hlt.bme.hu/semantics/external/pages/jokuc/en.wikipedia.org/wiki/Lakes_Plain_languages.html}
\end{itemize}

\section{Religious and Community Resources}
\begin{itemize}
\item Joshua Project, Aikwakai/Sikaritai people group profile: \url{https://joshuaproject.net/people_groups/10197}
\end{itemize}

\section{Digital and Audio Resources}
None were found.

\section{Curriculum Development Context}
No evidence of school use was found. The available material is limited to a general description and a lexical source (Clouse 1997). Curriculum development would require new documentation work with the community.
\end{document}